```latex
\documentclass{article}
\usepackage{amsmath,amssymb,amsthm,xcolor}
\input{_macros}

\newtheorem{lemma}{Lemma}
\newtheorem{theorem}{Theorem}

\begin{document}

\section*{Permitted external facts}

We shall use the following permitted external facts.

\begin{enumerate}
    \item \textbf{The Division Algorithm:} for any integers $a$ and $d$ with $d>0$, there exist unique integers $q$ and $r$ with
    \[
        a=dq+r
        \quad\text{and}\quad
        0\le r<d.
    \]
    Uniqueness implies that $r=0$ if and only if $d\mid a$, and $r\neq 0$ if and only if $d\nmid a$.

    \item \textbf{Euclid's Lemma:} if $p$ is prime and $p\mid ab$, then $p\mid a$ or $p\mid b$.

    \item \textbf{The Well-Ordering Principle:} every nonempty set of positive integers has a least element.

    \item \textbf{Closure:} the integers are closed under addition, subtraction, and multiplication.
\end{enumerate}

\begin{lemma}
For any integer $n$, if $7\mid n^2$, then $7\mid n$.
\end{lemma}

\begin{proof}
\textbf{Step 1. Primality of $7$.}

A positive integer is prime if it is greater than $1$ and has no divisors except $1$ and itself. By the Division Algorithm, for each $d\in\{2,3,4,5,6\}$, divide $7$ by $d$ to obtain remainder $r_d$:
\[
    7=2\cdot 3+1.
\]
Thus $r_2=1\neq 0$, hence $2\nmid 7$ by the Division Algorithm.

\[
    7=3\cdot 2+1.
\]
Thus $r_3=1\neq 0$, hence $3\nmid 7$ by the Division Algorithm.

\[
    7=4\cdot 1+3.
\]
Thus $r_4=3\neq 0$, hence $4\nmid 7$ by the Division Algorithm.

\[
    7=5\cdot 1+2.
\]
Thus $r_5=2\neq 0$, hence $5\nmid 7$ by the Division Algorithm.

\[
    7=6\cdot 1+1.
\]
Thus $r_6=1\neq 0$, hence $6\nmid 7$ by the Division Algorithm.

If $d>7$ and $d\mid 7$, then $7=dk$ for some positive integer $k$, giving
\[
    d=\frac{7}{k}\le \frac{7}{1}=7,
\]
a contradiction. The positivity of $k$ follows because $7>0$ and $d>0$, so $k=7/d>0$, and hence $k\ge 1$. If $d=7$, then $7\mid 7$ trivially. Thus the only positive divisors of $7$ are $1$ and $7$, so $7$ is prime.

\textbf{Step 2. Apply Euclid's Lemma.}

Suppose $7\mid n^2$. Since the integers are closed under multiplication,
\[
    n^2=n\cdot n
\]
is an integer. Since $7$ is prime by Step $1$ and $7\mid (n\cdot n)$, Euclid's Lemma gives $7\mid n$ or $7\mid n$; in either case, $7\mid n$.
\end{proof}

\begin{theorem}
$\sqrt{7}$ is irrational.
\end{theorem}

\begin{proof}
\textbf{Step 1. Setup.}

Assume for contradiction that $\sqrt{7}$ is rational. By definition, there exist integers $a,b$ with $b\neq 0$ such that
\[
    \sqrt{7}=\frac{a}{b}.
\]
Replacing $a$ by $-a$ and $b$ by $-b$ if necessary, noting that
\[
    \frac{-a}{-b}=\frac{a}{b},
\]
preserves the equality, so we may assume $b>0$. Since $\sqrt{7}$ is the positive square root of $7$ and $7>0$, we have $\sqrt{7}>0$. Multiplying both sides of
\[
    \sqrt{7}=\frac{a}{b}
\]
by $b>0$ gives
\[
    a=b\sqrt{7}.
\]
Since $b>0$ and $\sqrt{7}>0$, $a$ is a product of two positive numbers, so $a>0$.

\textbf{Step 2. Well-ordering.}

Therefore $b$ is a positive integer with
\[
    b\sqrt{7}=a\in\mathbb{Z}_{>0}.
\]
Define
\[
    S=\{q\in\mathbb{Z}_{>0}:q\sqrt{7}\in\mathbb{Z}_{>0}\}.
\]
Since $b\in S$, the set $S$ is nonempty. By the Well-Ordering Principle, $S$ has a least element $q_0$. Set
\[
    p_0=q_0\sqrt{7}\in\mathbb{Z}_{>0}.
\]
Then
\[
    \sqrt{7}=\frac{p_0}{q_0}.
\]

\textbf{Step 3. $7$ divides $p_0$.}

Squaring both sides of
\[
    \sqrt{7}=\frac{p_0}{q_0}
\]
gives
\[
    (\sqrt{7})^2=\left(\frac{p_0}{q_0}\right)^2,
\]
that is,
\[
    7=\frac{p_0^2}{q_0^2}.
\]
Multiplying both sides by $q_0^2$, which is a positive integer since $q_0>0$ and the integers are closed under multiplication, gives
\[
    p_0^2=7q_0^2.
\]
Hence $7\mid p_0^2$. By the Lemma, with $n=p_0$ and noting that $p_0\in\mathbb{Z}$, we get
\[
    7\mid p_0.
\]
Write
\[
    p_0=7k.
\]
Since $p_0>0$ and $7>0$, we have
\[
    k=\frac{p_0}{7}>0,
\]
so $k\in\mathbb{Z}_{>0}$.

\textbf{Step 4. $7$ divides $q_0$.}

Substituting $p_0=7k$ into
\[
    p_0^2=7q_0^2
\]
gives
\[
    (7k)^2=7q_0^2,
\]
so
\[
    49k^2=7q_0^2.
\]
Since $7\neq 0$, divide both sides by $7$ to obtain
\[
    7k^2=q_0^2.
\]
Also, $k^2=k\cdot k$ is a positive integer by closure. Hence
\[
    7\mid q_0^2.
\]
By the Lemma, with $n=q_0$ and noting that $q_0\in\mathbb{Z}$, we get
\[
    7\mid q_0.
\]
Write
\[
    q_0=7m.
\]
Since $q_0>0$ and $7>0$, we have
\[
    m=\frac{q_0}{7}>0,
\]
so $m\in\mathbb{Z}_{>0}$.

\textbf{Step 5. Contradiction.}

Now
\[
    \sqrt{7}=\frac{p_0}{q_0}=\frac{7k}{7m}.
\]
Since $m>0$, we have $m\neq 0$, so we may cancel the common factor $7$:
\[
    \sqrt{7}=\frac{k}{m}.
\]
Multiplying both sides by $m>0$ gives
\[
    m\sqrt{7}=k\in\mathbb{Z}_{>0}.
\]
Thus $m\in S$. But $q_0=7m$ and $7\ge 2>1$, so
\[
    m=\frac{q_0}{7}<q_0,
\]
since $7>1$ implies $1/7<1$, which implies
\[
    m=\frac{q_0}{7}<q_0\cdot 1=q_0.
\]
This contradicts $q_0$ being the least element of $S$. Therefore the assumption that $\sqrt{7}$ is rational is false. Hence $\sqrt{7}$ is irrational.
\end{proof}

\end{document}
```